\section{Limitations and conclusion}
\label{sec:conclusion}
\paragraph{Limitations.} \method{} decodes at roughly a token per second, which suits interactive use with short answers and offline jobs but not serving many users. It runs one sequence at a time, and long prompts take minutes. The compressed experts cost some agreement with the API on open-ended text, although retrieval is unaffected, and machines with more storage can keep the experts lossless. All results come from a single machine, with short prompts and short decode blocks, and we did not measure decode at long context. We also did not compare against other engines: the community projects we cite run related checkpoints with different quantization on other hardware, and porting them to this setting was beyond our scope. A Mac with a faster SSD or more memory would shift the balance between the techniques we measured. Finally, our screening study covers contexts up to 32K tokens and does not settle the question for much longer ones.

\paragraph{Conclusion.} We asked what expert parallelism becomes on a laptop whose only other device is an SSD. Reading experts straight into fixed GPU buffers replaces the token dispatch of a cluster, and an LRU cache replaces expert replication. A single decoding token turns each layer into a chain of computing, reading and computing again, and the SSD spends a third of every step waiting. The overlap techniques of a cluster recover little of that time, and naive prefetching even makes it worse, because its guesses delay the reads a token needs. Letting the guesses wait until those reads have landed is enough to fill the idle time, and it makes decode clearly faster without changing a single output token. With these pieces, a 24\,GB MacBook Air runs the 754B-parameter DeepSeek-V4.1-Flash at about 1.2 tokens per second, close to DeepSeek's own deployment. We release the engine, the converter, the tests and all measurement logs.
